\documentclass[a4paper,10.5pt]{article}
\usepackage[T1]{fontenc}\usepackage[utf8]{inputenc}
\usepackage[italian]{babel}
\usepackage[top=2.2cm,bottom=2.2cm,left=2.35cm,right=2.35cm]{geometry}
\usepackage{hyperref}\usepackage{xcolor}\definecolor{accent}{HTML}{1F4E78}\hypersetup{colorlinks=true,urlcolor=accent}
\pagestyle{empty}\setlength{\parindent}{0pt}\setlength{\parskip}{7pt}
\begin{document}
\textbf{Prabhavit Gupta}\\
Venezia, Italia \textbar{} +39 351 610 7807 \textbar{} prabhavitg@gmail.com\\[10pt]
Oggetto: Candidatura per Senior AI Engineer (rif. Lawrence Harvey)

Gentile team di selezione,

vi scrivo per candidarmi alla posizione di Senior AI Engineer che state ricercando per conto del vostro cliente. Ho tre anni di esperienza continuativa come sviluppatore AI e automazione presso otto aziende europee, con particolare focus su agenti AI e pipeline basate su LLM portate in produzione.

In FlowSpace ho costruito un bot AI che, a partire dalla brochure di un fornitore di spazi di lavoro in qualsiasi formato, redige l'intero annuncio con una persona che verifica il risultato, con oltre 900 bozze prodotte finora; ho inoltre realizzato agenti di supporto AI che rispondono alle richieste dei clienti leggendo dati live della piattaforma e passano a una persona le richieste a bassa confidenza, oltre a integrazioni API con OfficeRnD e Nexudus per le prenotazioni automatizzate.

In UltraStrategy ho costruito la pipeline di parsing dei CV integrata nella piattaforma: l'utente carica il CV, la pipeline estrae i dati e valuta la rilevanza del candidato per tipi di ruolo specifici, il risultato viene mostrato nel frontend e infine viene generato un "Ultra CV" standardizzato; per l'estrazione e la generazione ho usato Affinda e le API di Claude, e lo strumento di parsing è stato costruito con n8n ed Eden AI.

Vi segnalo con trasparenza che il mio percorso con i framework di orchestrazione agentica che citate, come LangGraph, LangChain e Google ADK, e con i protocolli MCP e A2A, è limitato: la mia esperienza pratica di orchestrazione è con n8n e con le API di Claude, incluso l'uso quotidiano di Claude Code come strumento di sviluppo assistito da AI. Ritengo che l'impianto concettuale sia lo stesso (agenti, tool, retrieval, valutazione) e sono pronto a colmare rapidamente il divario specifico sui framework indicati.

Vivo a Venezia, ho un permesso di soggiorno UE di lungo periodo e non necessito di sponsorizzazione per lavorare in Italia. Sono disponibile per la modalità ibrida richiesta dalla posizione, incluse le giornate in sede a Milano. Il mio italiano è conversazionale (B1-B2) e il mio inglese è professionale (C2). Sono disponibile con breve preavviso.

Grazie per l'attenzione. Resto a disposizione per un colloquio.

Cordiali saluti,\\
Prabhavit Gupta
\end{document}
